\documentclass[11pt]{article}
\usepackage[margin=1in]{geometry}
\usepackage{hyperref}
\usepackage{xcolor}
\usepackage{titlesec}
\usepackage{enumitem}
\usepackage{fancyhdr}
\usepackage{setspace}
\usepackage{parskip}

% Hyperlink styling
\hypersetup{
    colorlinks=true,
    linkcolor=blue,
    urlcolor=blue,
    citecolor=blue
}

% Section formatting
\titleformat{\section}{\Large\bfseries}{}{0em}{}
\titleformat{\subsection}{\large\bfseries}{}{0em}{}

% Header/Footer
\pagestyle{fancy}
\fancyhf{}
\fancyhead[L]{\textit{The Daily Gamecock}}
\fancyhead[R]{\textit{Free Speech Article -- Three Versions}}
\fancyfoot[C]{\thepage}
\renewcommand{\headrulewidth}{0.4pt}

\begin{document}

% Title Page
\begin{titlepage}
\centering
\vspace*{2cm}
{\Huge\bfseries Free Speech Hypocrisy at USC\par}
\vspace{0.5cm}
{\Large Three Article Versions\par}
\vspace{2cm}
{\large Prepared for The Daily Gamecock\par}
\vspace{0.5cm}
{\large January 2026\par}
\vspace{3cm}

\begin{tabular}{|l|p{8cm}|}
\hline
\textbf{Version} & \textbf{Description} \\
\hline
Version 1 & \textbf{Emotional/Passionate} -- Vivid imagery, rhetorical questions, appeals to outrage and moral clarity \\
\hline
Version 2 & \textbf{Fact-Driven/Analytical} -- Data-led, evidence accumulation, measured but argumentative \\
\hline
Version 3 & \textbf{Balanced/Medium} -- Mix of narrative and data, accessible, most likely to run \\
\hline
\end{tabular}

\vfill
{\small All versions include embedded hyperlinks to sources}
\end{titlepage}

\tableofcontents
\newpage

%%%%%%%%%%%%%%%%%%%%%%%%%%%%%%%%%%%%%%%%%%%%%%%%%%%%%%%%%%%%%
% VERSION 1: EMOTIONAL/PASSIONATE
%%%%%%%%%%%%%%%%%%%%%%%%%%%%%%%%%%%%%%%%%%%%%%%%%%%%%%%%%%%%%

\section{VERSION 1: Emotional/Passionate}
\vspace{1cm}

{\LARGE\bfseries Nobody at USC Actually Believes in Free Speech}

\vspace{0.3cm}
{\large\textit{The convenient principle everyone invokes and no one applies consistently has exposed the moral bankruptcy on every side of the campus culture war}}

\vspace{0.3cm}
By [AUTHOR NAME]

\vspace{0.5cm}
\hrule
\vspace{0.5cm}

\onehalfspacing

Picture this: September 18, 2024. A Wednesday evening at the University of South Carolina, and two Americas are unfolding simultaneously just a few hundred yards apart.

At Russell House, behind metal barricades and a wall of police officers, Milo Yiannopoulos and Gavin McInnes---the founder of the Proud Boys---are delivering their ``\href{https://www.dailygamecock.com/article/2024/09/over-1000-students-attend-blatt-bonanza-during-controversial-campus-event-news-grady}{Roast of Kamala Harris}'' to a sparse crowd. The room that could hold hundreds sits mostly empty. Outside, protesters gather, some holding signs, others just watching.

But walk across campus to Blatt Field, and you'll find a different scene entirely: \href{https://www.dailygamecock.com/article/2024/09/over-1000-students-attend-blatt-bonanza-during-controversial-campus-event-news-grady}{over 1,000 students} dancing, eating free food, taking selfies with President Michael Amiridis himself. ``Blatt Bonanza,'' they called it---a counter-event organized to draw attention away from the provocateurs in Russell House.

Both sides claimed victory for free speech that night. Uncensored America said they exercised their constitutional rights. The counter-programmers said they responded with ``more speech, not enforced silence''---the very language of Justice Brandeis himself.

And here's where I want to grab you by the shoulders and shake you: \textit{Did anyone there actually believe in the principle they were invoking?}

I don't think so. And I'm furious about it.

\subsection{The Conservative Fraud}

Let's start with the right, because their hypocrisy is currently the most dangerous.

Conservative students at USC demand---\textit{demand}---platforms for speakers like Myron Gaines to call women ``biologically inferior,'' for Gavin McInnes to spew his bile, for provocateurs whose entire business model is generating outrage. They wrap themselves in the First Amendment like it's a sacred garment. They invoke the marketplace of ideas. They accuse anyone who objects of being snowflakes who can't handle disagreement.

Fine. I can respect a principled commitment to free expression, even for speech I find repugnant.

But here's what makes my blood boil: these same defenders of unfettered discourse are cheering as state legislatures across the South systematically dismantle academic freedom. In Florida, Governor DeSantis's crusade has resulted in \href{https://www.thefire.org/news/worst-both-worlds-campus-free-speech}{13 DEI employees fired at the University of Florida}. At New College, 36 professors---a full third of the faculty---fled after a conservative board takeover turned their institution into an ideological battleground.

In Texas, SB 18 eliminates tenure protections entirely. SB 16 lets students get professors \textit{fired} for teaching about systemic racism. In our own state, South Carolina legislators passed a bill 84-30 in March 2024 prohibiting universities from considering political viewpoints in hiring---a law that sounds neutral until you realize it was designed to prevent hiring professors whose scholarship conservatives find objectionable.

Turning Point USA maintains a \href{https://www.thefire.org/research-learn/campus-deplatforming-database}{Professor Watchlist} targeting over 300 academics as ``Terror Supporters'' and ``Antifa.'' At Arizona State, TPUSA employees physically \textit{assaulted} Professor David Boyles. This isn't a free speech organization---it's an intimidation campaign.

So let me ask conservative students something directly: You cannot demand a platform for Milo Yiannopoulos while supporting laws that fire professors for their scholarship. You cannot cry censorship when a student government denies your funding request while applauding when state governments literally ban words from classrooms.

That's not free speech. That's \textit{your} speech. There's a difference, and pretending otherwise is a lie.

\subsection{The Progressive Fraud}

But don't you dare think I'm letting the left off the hook. Oh no. Progressives have their own cathedral of hypocrisy, and it's equally infuriating.

Consider Stanford Law School in 2023, where students shouted down Judge Kyle Duncan for his entire scheduled appearance. The DEI Dean confronted him, asking if ``the juice is worth the squeeze''---a barely veiled threat wrapped in academic jargon. After that incident, \href{https://www.thefire.org/news/worst-both-worlds-campus-free-speech}{conservative student comfort dropped from 45\% to 6\%}. Six percent. At an elite law school. In America.

Or Middlebury in 2017, where Professor Allison Stanger---a \textit{liberal} professor who was merely \textit{moderating} a talk by Charles Murray---was \href{https://www.thefire.org/research-learn/campus-deplatforming-database}{hospitalized with a concussion} after protesters attacked her. She later wrote that she feared for her life. For moderating a conversation.

Here at USC, \href{https://www.postandcourier.com/education-lab/usc-doj-charlie-kirk-outdoor-event-controversial/article_c1d6ca0b-1173-403d-bd16-6f8fec66f748.html}{27,000 people signed a petition} demanding the McInnes and Yiannopoulos event be cancelled. Twenty-seven thousand signatures to silence speakers. And these same signatories would be the first to defend their right to protest, to occupy buildings, to disrupt campus operations when \textit{their} causes are at stake.

In April 2024, pro-Palestinian protesters at USC were arrested for ``breach of peace.'' Many of those same students---or their ideological allies---would argue that disruptive protest is protected speech when it serves progressive goals, but that conservative speakers shouldn't be given platforms at all.

According to \href{https://www.thefire.org/research-learn/campus-deplatforming-database}{FIRE's data}, heckler's veto incidents---where speakers are shouted down or physically prevented from speaking---skyrocketed from 3 in 2021 to 79 in 2024. The \href{https://knightfoundation.org/reports/college-student-views-on-free-expression-and-campus-speech-2024/}{Knight Foundation found} that 40\% of college students now consider shouting down speakers ``sometimes'' or ``always'' acceptable. Thirty-three percent---up from 20\% in 2020---consider \textit{violence} sometimes acceptable to stop a campus speech.

Here's the contradiction that should haunt every progressive who claims to believe in free expression: You cannot argue that shouting down speakers is ``counter-speech'' protected by the First Amendment while \textit{also} arguing that speakers you dislike shouldn't receive platforms in the first place. Either disruptive responses to speech are legitimate or they aren't. Pick one.

\subsection{The Administrative Fraud}

And then there's the university itself, playing both sides with the cynicism of a career politician.

President Amiridis says USC follows the Chicago Principles. He's \href{https://sc.edu/about/our_leadership/president/presidential-communications/2024-08-27_statement_about_campus_speakers.php}{on record stating} that ``the solution to offensive speech is more speech.'' Beautiful sentiment. I'd applaud it if the institution actually lived by it.

But here's what ``more speech'' looked like on September 18, 2024: A university-supported counter-event with free food, entertainment, and the president himself taking selfies, while controversial speakers addressed empty rooms behind security barricades. The institution took a side---it just won't admit it.

The numbers are damning. At UC Berkeley, security for Ben Shapiro cost \href{https://www.berkeleyside.org/2018/02/06/uc-berkeley-spent-close-4m-security-just-one-month-2017}{\$836,421}. For Milo Yiannopoulos? \$800,000 for a fifteen-minute appearance. For Justice Sonia Sotomayor? Five thousand dollars.

The Supreme Court ruled in \href{https://supreme.justia.com/cases/federal/us/505/123/}{\textit{Forsyth County v. Nationalist Movement}} that this exact practice is unconstitutional: ``Speech cannot be financially burdened, any more than it can be punished or banned, simply because it might offend a hostile mob.'' When universities charge student organizations security fees based on how controversial their speakers are, they're letting the mob dictate who gets to speak. That's not neutrality. That's capitulation dressed in bureaucratic language.

In September 2025, the Department of Justice sent USC a letter citing ``\href{https://www.postandcourier.com/education-lab/usc-doj-charlie-kirk-outdoor-event-controversial/article_c1d6ca0b-1173-403d-bd16-6f8fec66f748.html}{troubling allegations}'' that the university told conservative groups like Turning Point USA and Uncensored America that they couldn't hold outdoor events for ``controversial speakers.'' The Student Government voted to reject Uncensored America's \$3,600 funding request---a content-based decision by any honest assessment.

So much for viewpoint neutrality.

\subsection{The Numbers Don't Lie}

Here's what the data shows, and it should make everyone uncomfortable.

According to \href{https://www.thefire.org/research-learn/campus-deplatforming-database}{FIRE's deplatforming database}, 1,406 deplatforming attempts have occurred since 1998, with 39\% succeeding. From 1998 to 2012, 72\% of those attempts came from the \textit{right}---the Cardinal Newman Society targeting ``The Vagina Monologues,'' conservative Catholics protesting Obama's commencement speech at Notre Dame. From 2013 to 2022, the script flipped: 70\% came from the left. Since 2023, it's been roughly equal.

In 2025, government censorship exploded: \href{https://www.thefire.org/news/worst-both-worlds-campus-free-speech}{114 incidents involving politicians} compared to 102 total incidents from 2000 to 2024 combined.

The pattern is clear: Whichever side has institutional power uses it to silence the other. The principle of ``free speech'' is invoked exclusively when convenient and abandoned the moment it becomes inconvenient.

USC's own trajectory proves the point. In 2024, FIRE ranked USC \href{https://thefire.org/sites/default/files/2024-09/2025\%20CFSR\%20Spotlight\%20-\%20U\%20of\%20South\%20Carolina.pdf}{246th out of 251 schools} for free speech---bottom five in the country. After adopting the Chicago Principles, we climbed to 34th in 2025, then 22nd in 2026.

Improvement came from \textit{policy reform}, not from either side suddenly discovering principles they'd ignored for years.

\subsection{The Only Honest Positions}

There are two intellectually honest positions on campus speech. Just two.

The first is consistent viewpoint neutrality: USC provides equal institutional support for all legal speech regardless of content. Security costs are never passed to student organizations. No funding decisions are based on speaker ideology. Protest is protected; disruption that prevents speech is not. Apply this rule to everyone, always, with no exceptions.

The second is openly acknowledged value judgments: USC admits it makes content-based decisions about which speech to amplify and debates those values transparently. Counter-programming for speakers deemed hateful becomes official policy. The community decides collectively what belongs on our campus and what doesn't.

Both positions are defensible. What's not defensible---what's intellectually cowardly and morally bankrupt---is what we have now: pretending to neutrality while taking sides, invoking principles we apply selectively, and calling our opponents hypocrites while being hypocrites ourselves.

\subsection{The Question That Exposes Everything}

The next time you hear someone at USC invoke ``free speech''---whether it's Uncensored America demanding a platform, protesters demanding a cancellation, or administrators claiming neutrality---ask them one question:

\textit{Would you apply this principle if the speaker's politics were reversed?}

If a conservative speaker deserves a platform, does a Marxist professor deserve tenure protections? If progressive students can shout down a judge, can conservative students shout down a DEI speaker? If the university can spend resources on counter-programming against right-wing provocateurs, can it do the same against left-wing activists?

If the answer to any of these is ``no''---if the principle only applies when it benefits your side---then you don't believe in free speech.

You believe in power.

And at least have the courage to admit it.

\vspace{0.5cm}
\hrule
\vspace{0.3cm}
{\small\textit{The opinions expressed in this article are those of the author and do not necessarily reflect the views of The Daily Gamecock.}}

\newpage

%%%%%%%%%%%%%%%%%%%%%%%%%%%%%%%%%%%%%%%%%%%%%%%%%%%%%%%%%%%%%
% VERSION 2: FACT-DRIVEN/ANALYTICAL
%%%%%%%%%%%%%%%%%%%%%%%%%%%%%%%%%%%%%%%%%%%%%%%%%%%%%%%%%%%%%

\section{VERSION 2: Fact-Driven/Analytical}
\vspace{1cm}

{\LARGE\bfseries The Free Speech Fraud at USC}

\vspace{0.3cm}
{\large\textit{Everyone claims to defend the First Amendment. The data shows they only defend it for themselves.}}

\vspace{0.3cm}
By [AUTHOR NAME]

\vspace{0.5cm}
\hrule
\vspace{0.5cm}

\onehalfspacing

The numbers tell a story that no one at the University of South Carolina wants to hear. Since 1998, there have been \href{https://www.thefire.org/research-learn/campus-deplatforming-database}{1,406 documented attempts} to deplatform speakers at American universities, with 39 percent succeeding. From 1998 to 2012, 72 percent of those attempts came from the political right. From 2013 to 2022, 70 percent came from the left. Today, the attempts split roughly evenly between both sides.

The pattern is unmistakable: whichever faction holds institutional power uses it to silence the other. ``Free speech'' is not a principle at USC or anywhere else. It is a weapon, wielded by whoever finds it momentarily convenient and discarded the instant it becomes inconvenient.

Consider September 18, 2024. Two events unfolded simultaneously on our campus. At Russell House, Milo Yiannopoulos and Gavin McInnes addressed a sparse audience behind metal barricades, their ``Roast of Kamala Harris'' drawing more police officers than attendees. Across campus at Blatt Field, over \href{https://www.dailygamecock.com/article/2024/09/over-1000-students-attend-blatt-bonanza-during-controversial-campus-event-news-grady}{1,000 students} danced, ate free food, and took selfies with President Michael Amiridis at ``Blatt Bonanza,'' a university-supported counter-event.

Both sides claimed victory for free speech. Uncensored America said they exercised their constitutional rights despite opposition. Counter-programmers argued they responded with ``more speech, not enforced silence,'' invoking Justice Brandeis's famous formulation from \href{https://supreme.justia.com/cases/federal/us/274/357/}{\textit{Whitney v. California}}. But the data suggests neither side believed in the principle they invoked. They believed in winning.

\subsection{The Conservative Contradiction}

Conservative students at USC demand platforms for provocative speakers while supporting policies that systematically restrict academic speech elsewhere. The cognitive dissonance is quantifiable.

Turning Point USA maintains a ``Professor Watchlist'' targeting over 300 academics as supposed threats. At Arizona State University, the consequences of such targeting turned violent when TPUSA employees \href{https://www.insidehighered.com/news/2017/12/19/arizona-state-student-events-official-assaulted-political-group}{physically assaulted} Professor David Boyles. The organization that demands campus platforms simultaneously encourages the harassment of professors whose speech they oppose.

The legislative record is equally damning. In Florida, Governor Ron DeSantis's crusade against ``woke ideology'' resulted in \href{https://pen.org/report/educational-gag-orders/}{13 DEI employees} fired at the University of Florida. At New College, 36 professors---one-third of the faculty---fled after a conservative board takeover restructured the institution's academic mission. Texas's SB 18 eliminates tenure protections entirely, while SB 16 empowers students to initiate termination proceedings against professors who teach concepts related to critical race theory.

South Carolina has not been immune. In March 2024, state legislators \href{https://www.scstatehouse.gov/sess125_2023-2024/bills/3728.htm}{passed a bill} prohibiting universities from factoring political stances into hiring decisions by an 84-30 vote---a content-based restriction on institutional speech dressed up as viewpoint neutrality.

As \href{https://pen.org/report/educational-gag-orders/}{PEN America} documented: ``There is no act of censorship that is more explicit than literally banning words.'' You cannot coherently demand a platform for speakers who call women ``biologically inferior'' while supporting legislation that fires professors for teaching about systemic racism. That is not a free speech position. That is a ``my speech'' position.

\subsection{The Progressive Contradiction}

Progressive students fare no better under scrutiny. The data on campus disruptions reveals a movement that claims ``deplatforming isn't censorship'' while demanding their own protests remain absolutely protected.

The heckler's veto---organized disruption intended to prevent speakers from being heard---has exploded. According to FIRE, incidents rose from \href{https://www.thefire.org/news/worst-both-worlds-campus-free-speech}{3 in 2021 to 79 in 2024}. At Stanford Law School in 2023, students shouted down Judge Kyle Duncan for the entirety of his scheduled appearance. The DEI dean confronted him mid-event, asking whether ``the juice was worth the squeeze.'' Conservative students' reported comfort on campus \href{https://www.thefire.org/news/stanford-students-comfort-expressing-ideas-dropped-6-following-judge-duncan-event}{plummeted from 45 percent to 6 percent}.

The incidents have turned violent. At Middlebury College in 2017, Professor Allison Stanger was \href{https://www.nytimes.com/2017/03/13/opinion/understanding-the-angry-mob-that-gave-me-a-concussion.html}{hospitalized with a concussion} while moderating an event with Charles Murray. ``I feared for my life,'' she later wrote.

At USC specifically, \href{https://www.postandcourier.com/education-lab/usc-doj-charlie-kirk-outdoor-event-controversial/article_c1d6ca0b-1173-403d-bd16-6f8fec66f748.html}{27,000 people signed a petition} to cancel the McInnes-Yiannopoulos event---the same community that would vigorously defend students' right to protest. In April 2024, pro-Palestinian demonstrators were arrested for breach of peace, their supporters decrying the suppression of protected speech. But those same supporters had signed petitions to prevent speech they opposed.

The logical contradiction is inescapable. Either disruptive response to speech is legitimate counter-expression, or it is an impermissible silencing tactic. You cannot claim both positions depending on whose ox is being gored.

\subsection{The Administrative Contradiction}

Perhaps the most consequential hypocrisy belongs to university administrators, who claim viewpoint neutrality while making content-based decisions that determine whose speech gets heard.

The security cost mechanism is the clearest example. UC Berkeley spent \href{https://www.berkeleyside.org/2018/02/06/uc-berkeley-spent-close-4m-security-just-one-month-2017}{\$836,421 on security} for Ben Shapiro's 2017 appearance. For a 15-minute Milo Yiannopoulos speech, the tab reached \$800,000. Justice Sonia Sotomayor's 2011 visit cost approximately \$5,000.

The Supreme Court addressed this disparity directly in \href{https://supreme.justia.com/cases/federal/us/505/123/}{\textit{Forsyth County v. Nationalist Movement}} (1992): ``Speech cannot be financially burdened, any more than it can be punished or banned, simply because it might offend a hostile mob.'' Charging speakers more because their content provokes opposition is unconstitutional. Yet universities routinely pass security costs to student organizations, creating a de facto system where controversial speech becomes unaffordable.

USC is now under federal scrutiny for precisely this issue. In September 2025, the Department of Justice sent a letter citing ``\href{https://www.postandcourier.com/education-lab/usc-doj-charlie-kirk-outdoor-event-controversial/article_c1d6ca0b-1173-403d-bd16-6f8fec66f748.html}{troubling allegations}'' that the university told Turning Point USA and Uncensored America that no outdoor events would be permitted for ``controversial speakers.'' That is a textbook content-based restriction---the government deciding which viewpoints warrant public space.

Student Government compounds the problem. When SG voted to reject Uncensored America's \$3,600 funding request, it made a decision based on the content of proposed speech. Meanwhile, the administration funded Blatt Bonanza as counter-programming---a choice that, however defensible as ``more speech,'' represents an institutional thumb on the scale.

President Amiridis \href{https://sc.edu/about/our_leadership/president/presidential-communications/2024-08-27_statement_about_campus_speakers.php}{stated publicly} that USC follows the Chicago Principles and that ``the solution to offensive speech is more speech.'' But when that ``more speech'' means 1,000 students at a university-supported celebration while controversial speakers address empty rooms behind barricades, the institution has taken a side. Claiming otherwise is administrative doublespeak.

\subsection{The Survey Data Nobody Wants}

The \href{https://knightfoundation.org/reports/college-student-views-on-free-expression-and-campus-speech-2024/}{Knight Foundation's 2024 survey} quantifies the erosion of free speech as a shared value. Forty percent of college students now consider shouting down speakers ``sometimes'' or ``always'' acceptable. More alarmingly, 33 percent consider violence to prevent speech acceptable---up from 20 percent in 2020.

The self-censorship figures are equally stark: 66 percent of students report self-censoring in the classroom. When two-thirds of a university community feels unable to speak freely, something has broken in the culture of open inquiry.

USC's trajectory reflects these national patterns. In 2024, \href{https://thefire.org/sites/default/files/2024-09/2025\%20CFSR\%20Spotlight\%20-\%20U\%20of\%20South\%20Carolina.pdf}{FIRE ranked USC} 246th out of 251 institutions---bottom five in the nation. After adopting the Chicago Principles, the ranking improved to 34th in 2025 and 22nd in 2026. The improvement came from policy reform, not from students or faculty suddenly believing in the principle. Institutional structure changed behavior. Conviction had nothing to do with it.

\subsection{Two Honest Positions}

Intellectual honesty requires acknowledging that only two coherent positions exist on campus speech.

The first is consistent viewpoint neutrality: the university provides equal institutional support for all legal speech regardless of content. Security costs are never passed to student organizations. No funding decisions reference speaker ideology. Protest is protected; disruption that prevents speech is sanctioned. The institution genuinely does not care whether speakers are progressive, conservative, or anything else.

The second is openly acknowledged value judgment: the university admits it makes content-based decisions about which speech to amplify and debates those values transparently. Counter-programming for speakers deemed hateful becomes stated institutional policy. The community decides collectively, through legitimate governance, what belongs on its campus.

What USC has now is the worst option: pretending to neutrality while taking sides, invoking principles applied selectively, and calling opponents hypocrites while practicing hypocrisy ourselves.

The next time someone at USC invokes ``free speech''---whether it is Uncensored America demanding a platform, protesters demanding a cancellation, or administrators claiming neutrality---ask one question: Would you apply this principle if the speaker's politics were reversed?

If the answer is no, they do not believe in free speech. They believe in power. And until we admit that honestly, every debate about campus speech will remain an exercise in bad faith, dressed up in constitutional language that nobody actually means.

\vspace{0.5cm}
\hrule
\vspace{0.3cm}
{\small\textit{The Daily Gamecock is the independent student newspaper of the University of South Carolina.}}

\newpage

%%%%%%%%%%%%%%%%%%%%%%%%%%%%%%%%%%%%%%%%%%%%%%%%%%%%%%%%%%%%%
% VERSION 3: BALANCED/MEDIUM
%%%%%%%%%%%%%%%%%%%%%%%%%%%%%%%%%%%%%%%%%%%%%%%%%%%%%%%%%%%%%

\section{VERSION 3: Balanced/Medium (Recommended)}
\vspace{1cm}

{\LARGE\bfseries Nobody at USC Actually Believes in Free Speech}

\vspace{0.3cm}
{\large\textit{Everyone claims the principle. Nobody applies it consistently. It's time we admitted we're all hypocrites.}}

\vspace{0.3cm}
By [AUTHOR NAME]

\vspace{0.5cm}
\hrule
\vspace{0.5cm}

\onehalfspacing

September 18, 2024. Two events unfolded simultaneously on USC's campus, and both sides walked away claiming victory for free speech.

At Russell House, Milo Yiannopoulos and Proud Boys founder Gavin McInnes addressed a sparse crowd from behind security barricades for their ``Roast of Kamala Harris.'' Across campus at Blatt Field, \href{https://www.dailygamecock.com/article/2024/09/over-1000-students-attend-blatt-bonanza-during-controversial-campus-event-news-grady}{over 1,000 students} danced, ate free food, and took selfies with President Amiridis at ``Blatt Bonanza,'' a counter-programming event organized to drain attention from the controversial speakers.

Uncensored America celebrated exercising their First Amendment rights. Counter-programmers celebrated responding with ``more speech, not enforced silence.'' Everyone invoked the same principle. Everyone claimed the moral high ground.

But here's the question nobody wanted to ask that night: Did anyone on either side actually believe in free speech as a principle, or were they just wielding it as a weapon?

After months of research into campus speech controversies at USC and nationwide, I've reached an uncomfortable conclusion: Nobody at this university genuinely supports free speech. What they support is \textit{their} speech. And until we're honest about that, every argument about speakers, protests, and censorship is theater.

\subsection{The Conservative Hypocrisy}

Let's start with the groups demanding platforms for provocative speakers. Conservative student organizations at USC and across the country position themselves as free speech warriors, fighting against liberal censorship in academia. And they have legitimate grievances---\href{https://www.thefire.org/research-learn/campus-deplatforming-database}{FIRE's database} documents 1,406 deplatforming attempts since 1998, with the majority in recent years coming from the left.

But the same movement demanding stages for Milo Yiannopoulos supports state policies that fire professors for their viewpoints.

In Florida, Governor DeSantis's crusade against ``woke'' education resulted in 13 DEI employees terminated at the University of Florida. At New College, a conservative board takeover drove 36 professors---one-third of the faculty---to flee. Texas eliminated tenure protections and passed legislation allowing students to get professors fired for teaching critical race theory. And here in South Carolina, a March 2024 bill prohibiting universities from considering political stances in hiring \href{https://www.scstatehouse.gov/}{passed 84-30}.

Turning Point USA, one of the loudest voices demanding campus platforms, maintains a ``Professor Watchlist'' labeling over 300 academics as ``Terror Supporters'' and ``Antifa.'' At Arizona State, a professor on that list was physically assaulted by TPUSA employees.

Here's the contradiction: You cannot demand a platform for speakers who call women ``biologically inferior'' while simultaneously supporting legislation that fires professors for teaching about systemic racism. You cannot invoke free speech to defend provocateurs while backing government efforts to ban certain words and ideas from classrooms. That's not free speech---that's \textit{your} speech.

As \href{https://pen.org/}{PEN America} noted in their analysis of educational gag orders: ``There is no act of censorship that is more explicit than literally banning words.''

\subsection{The Progressive Hypocrisy}

But before progressive readers feel vindicated, let's examine the other side of the ledger.

The argument that ``deplatforming isn't censorship'' has become common on the left. Shouting down speakers, the logic goes, is itself protected speech---a form of counter-speech that exercises, rather than suppresses, First Amendment rights. Yet the same students who defend disruptive protest demand protection for their own demonstrations.

At Stanford Law School in 2023, students shouted down federal judge Kyle Duncan for his entire scheduled appearance. The school's DEI dean confronted him publicly, asking whether ``the juice is worth the squeeze.'' According to FIRE surveys, conservative student comfort at Stanford plummeted from 45\% to 6\% after the incident.

The consequences of this approach can turn violent. At Middlebury College in 2017, Professor Allison Stanger was hospitalized with a concussion after protesters attacked her for merely moderating a conversation with Charles Murray. ``I feared for my life,'' she later wrote.

Here at USC, \href{https://www.postandcourier.com/education-lab/usc-doj-charlie-kirk-outdoor-event-controversial/article_c1d6ca0b-1173-403d-bd16-6f8fec66f748.html}{27,000 people signed a petition} demanding the cancellation of the McInnes/Yiannopoulos event. In April 2024, pro-Palestinian protesters were arrested for ``breach of peace''---many of them the same students who would vigorously defend their right to protest.

The contradiction: You cannot argue that shouting down speakers is ``counter-speech'' protected by the First Amendment while also arguing that speakers you dislike shouldn't be given platforms in the first place. Either disruptive response to speech is legitimate, or it isn't. You don't get to decide based on whether you agree with the target.

Justice Brandeis wrote in 1927 that ``the remedy to be applied is more speech, not enforced silence.'' The Blatt Bonanza organizers understood this. But a petition demanding cancellation isn't more speech---it's a demand for enforced silence with extra steps.

\subsection{The Administrative Hypocrisy}

Perhaps no one plays this game more cynically than university administrators, who claim viewpoint neutrality while making obviously content-based decisions.

Consider security costs. When Ben Shapiro spoke at UC Berkeley in 2017, the university spent \href{https://www.berkeleyside.org/2018/02/06/uc-berkeley-spent-close-4m-security-just-one-month-2017}{\$836,421 on security}. When Justice Sonia Sotomayor visited in 2011, security cost \$5,000. Milo Yiannopoulos's 15-minute Berkeley appearance required \$800,000 in protection.

The Supreme Court addressed this disparity directly in \href{https://supreme.justia.com/cases/federal/us/505/123/}{\textit{Forsyth County v. Nationalist Movement}} (1992): ``Speech cannot be financially burdened, any more than it can be punished or banned, simply because it might offend a hostile mob.'' Charging speakers or organizations more because their ideas provoke opposition is unconstitutional. Yet universities do it routinely.

USC is not immune. In September 2025, the Department of Justice sent the university a letter citing ``\href{https://www.postandcourier.com/education-lab/usc-doj-charlie-kirk-outdoor-event-controversial/article_c1d6ca0b-1173-403d-bd16-6f8fec66f748.html}{troubling allegations}'' about viewpoint discrimination. According to the DOJ, USC allegedly told conservative groups like TPUSA and Uncensored America that they couldn't hold outdoor events for ``controversial speakers''---a content-based restriction that, if true, violates the First Amendment at a public university.

Meanwhile, Student Government voted to reject Uncensored America's \$3,600 funding request---another content-based decision. And while students debate who pays for security at conservative events, nobody has asked how much USC spent on Blatt Bonanza, the administration-endorsed counter-event.

President Amiridis claims USC follows the Chicago Principles and that ``the solution to offensive speech is more speech.'' That sounds principled. But when ``more speech'' means 1,000 students at a university-supported party while controversial speakers address empty rooms behind barricades, the institution has taken a side. It just won't admit it.

\subsection{The Data Nobody Wants to Hear}

Here's where both sides will hate me equally.

According to \href{https://www.thefire.org/news/worst-both-worlds-campus-free-speech}{FIRE's research}, from 1998 to 2012, 72\% of campus deplatforming attempts came from the \textit{right}. The Cardinal Newman Society targeted ``The Vagina Monologues'' at Catholic schools. Conservative groups protested Obama and Clinton commencement speeches. The right was the censorship threat.

Then something shifted. From 2013 to 2022, 70\% of deplatforming attempts came from the \textit{left}. Trigger warnings, speaker cancellations, bias response teams---progressives controlled campus culture and used that power accordingly.

Since 2023, the attempts have been roughly equal. And in 2025, government censorship surged---FIRE documented 114 incidents involving politicians compared to just 102 total from 2000 to 2024 combined.

The pattern is clear: Whichever side has institutional power uses it to silence the other. The principle of ``free speech'' is invoked only when convenient.

The \href{https://knightfoundation.org/reports/college-student-views-on-free-expression-and-campus-speech-2024/}{Knight Foundation's 2024 survey} revealed the hypocrisy in stark numbers. Forty percent of students consider shouting down speakers ``sometimes'' or ``always'' acceptable. Thirty-three percent---up from 20\% in 2020---consider violence acceptable in some circumstances. Sixty-six percent report self-censoring in the classroom.

Everyone claims to support free speech. The data says otherwise.

\subsection{What Intellectual Honesty Would Require}

Only two honest positions exist on campus speech.

The first is consistent viewpoint neutrality: USC provides equal institutional support for all legal speech, regardless of content. Security costs are never passed to student organizations. No funding decisions are based on speaker ideology. Protest is protected; disruption that prevents speech is not. This is hard. It means defending the speech rights of people you find abhorrent. It means treating Milo Yiannopoulos and Angela Davis identically under policy.

The second is openly acknowledged value judgments: USC admits it makes content-based decisions about which speech to amplify and debates those values openly. Counter-programming for hateful speakers becomes official institutional policy. The community decides collectively what belongs on campus. This is also hard. It means abandoning the pretense of neutrality and accepting responsibility for those choices.

What we have now is the worst of both worlds---pretending to neutrality while taking sides, invoking principles we apply selectively, and calling opponents hypocrites while being hypocrites ourselves.

USC's FIRE ranking has improved dramatically---from \href{https://thefire.org/sites/default/files/2024-09/2025\%20CFSR\%20Spotlight\%20-\%20U\%20of\%20South\%20Carolina.pdf}{bottom five in 2024} to 22nd in 2026---largely because of policy reforms like adopting the Chicago Principles. That improvement came from institutional changes, not from either side suddenly discovering principled commitment to free expression.

So here's my challenge: The next time you hear someone at USC invoke ``free speech''---whether it's Uncensored America demanding a platform, protesters demanding a cancellation, or administrators claiming neutrality---ask one question.

\textit{Would you apply this principle if the speaker's politics were reversed?}

If the answer is no, they don't believe in free speech.

They believe in power.

And so, probably, do you.

\vspace{0.5cm}
\hrule
\vspace{0.3cm}
{\small\textit{The opinions expressed in this article are those of the author and do not necessarily reflect the views of The Daily Gamecock.}}

\end{document}
